% Standalone source. Intentionally not compiled.
% All theorem proofs needed for the stated mathematical claims are included.
% The companion Lean verification concerns only the explicitly identified algebraic core.
\documentclass[11pt]{article}
\usepackage[T1]{fontenc}
\usepackage[utf8]{inputenc}
\usepackage[margin=1in]{geometry}
\usepackage{amsmath,amssymb,amsthm,mathtools,enumitem}
\usepackage[hidelinks]{hyperref}
\newtheorem{theorem}{Theorem}[section]
\newtheorem{proposition}[theorem]{Proposition}
\newtheorem{lemma}[theorem]{Lemma}
\newtheorem{corollary}[theorem]{Corollary}
\theoremstyle{definition}
\newtheorem{definition}[theorem]{Definition}
\newtheorem{example}[theorem]{Example}
\theoremstyle{remark}
\newtheorem{remark}[theorem]{Remark}
\newcommand{\R}{\mathbb R}
\newcommand{\Set}{\mathbf{Set}}
\newcommand{\FinSet}{\mathbf{FinSet}}
\newcommand{\Mat}{\mathsf{Mat}}
\newcommand{\Kar}{\mathsf{Kar}}
\newcommand{\im}{\operatorname{im}}
\newcommand{\tr}{\operatorname{tr}}
\newcommand{\rank}{\operatorname{rank}}
\newcommand{\Hom}{\operatorname{Hom}}
\newcommand{\id}{\operatorname{id}}
\newcommand{\supp}{\operatorname{supp}}
\newcommand{\Stab}{\operatorname{Stab}}
\newcommand{\End}{\operatorname{End}}
\newcommand{\Bool}{\mathsf{Bool}}
\newcommand{\C}{\mathcal C}
\newcommand{\N}{\mathcal N}
\newcommand{\Fr}{\operatorname{Fr}}
\title{Operational Invention from Interaction\\
\large Constructed roles, acquired references, and resource-relative capability}
\author{}\date{5 November 2026}
\begin{document}
\maketitle
\begin{abstract}
An interaction can determine a reusable role, and executing it can produce a retained reference for later operations. These are different formation steps. This note proves a staged mechanism connecting them while distinguishing a semantic description, a native implementation, and the physically available resource state.

For a certified constant-rank response, finite coefficient arithmetic constructs its projective image role and preparation operations. A global scalar or free-coordinate interface need not exist; a finite sign-cycle criterion decides this obstruction for a certified rank-one class. Gaussian excitation with bounded-input rejection then acquires an actual well-conditioned frame of the response space at the retained operating point, almost surely in finite time with a uniform expected-cost bound. The acquired frame supports exact compositional coordinates and subsequent identification and controller synthesis under explicit finite-dimensional response promises.

The general categorical link is retained-context use as pullback to a token base. In group-action semantics, a retained $G/H$ reference enables exactly the $H$-equivariant operations; intersecting stabilizers describes cumulative acquisition. A separate finite native compiler establishes strict experimental gain on unchanged boundary sorts under a specified no-reference resource contract. This does not identify the matrix and group languages or infer physical realizability from equivariance alone.

The results establish constructive role formation and resource-relative capability, not unrestricted invention of new physical operations at unchanged access. Every complete acquisition-and-use procedure remains expressible in its stipulated runtime. The stronger endogenous-invention target is therefore not claimed solved. The categorical, algebraic, topological, and expressiveness ingredients have substantial prior art; no historical novelty claim is made.
\end{abstract}
\tableofcontents

\section{The operational question and the exact answer}
The question is not merely whether an interaction can be assigned a new symbol. It is whether a source interaction can produce something that has a warranted role, can be passed to later interactions, and changes what an investigator can subsequently do. A useful answer must specify the old language, the actual source access, the acquired resource, the admissible continuations, and the comparison used to call the result new.

We separate four levels throughout:
\begin{enumerate}
\item A \emph{semantic role} specifies how values participate in admitted preparations, tests, and transformations.
\item A \emph{native implementation} gives actual programs, representations, and adapters realizing that role.
\item A \emph{resource state} records the samples, devices, frames, storage, and operating conditions physically available to the current client.
\item A \emph{reuse theorem} states which later compositions preserve the role's guarantees, with the same retained resources.
\end{enumerate}
The same semantic role may have several implementations and resource costs. A formal image, colimit, or representable presheaf supplies none of those physical facts by itself.

The central construction is one chain. A response matrix $H$ determines a native image projector $e$. Its image is introduced as a first-class role $P_H$, with preparation and elimination operations. A frame of this role may fail to exist globally over operating conditions. Executing $H$ on bounded accepted random probes nevertheless produces a frame over the \emph{actual retained condition}. The token is then kept on every use, allowing coordinate operations and further experiments. The frame does not erase the global obstruction: it changes the base of operation from a condition $x$ to a condition together with an acquired reference $(x,V)$.

The finite symmetry model provides an exact resource-relative expressiveness calculation for the same retained-context mechanism. It is a controlled second realization, not an assertion that all geometric observations are invariant Boolean tests. This distinction matters: ambient coordinate readouts in the matrix model can already distinguish vectors that a name-blind symmetry model cannot.

\begin{theorem}[The closure boundary]\label{thm:closure}
Let an old program repertoire be closed under its declared composition rules. Suppose every new interface has an old representation, every extended primitive has a uniform old implementation, and these implementations preserve the common boundary observations and compose. Include acquisition, initialization, memory, and all later uses in these implementations. Then every finite extended experiment with old boundary types has an old implementation with the same observations. The same conclusion for recursive or unbounded experiments requires preservation of the corresponding recursion or iteration semantics. Resource-bounded conclusions additionally require preservation of the relevant costs.
\end{theorem}
\begin{proof}
Replace new primitives by their old implementations and passages through new types by their representation adapters. Induction on finite program formation proves equality of boundary observations. Compose the acquisition implementation with the compiled continuation. If an iteration or recursion constructor occurs, apply its assumed semantic preservation rather than finite induction alone. The cost statement uses the additional cost-preservation premise; extensional equality does not establish it.
\end{proof}

Thus a strict gain must concern a declared resource state, a restricted type grammar, a cost or uniformity requirement, or genuinely additional access. The theorem does not deny discoveries, useful abstractions, or fabrication. It prevents claiming that a closed implementation language transcends its own extensional image merely by running one of its programs. Felleisen's contextual-equivalence obstruction to macro elimination is a related established obstruction under its own translation conditions \cite[Theorem 3.14]{Felleisen1991}.

\begin{proposition}[Evidence and native feasibility are separate]\label{prop:evidence}
Let $q:W\to E$ send native worlds to all supplied evidence, and let $A_w$ be the implementations certified to realize a requested role in world $w$. Any sound evidence-only selection at $z$ must belong to $\bigcap_{q(w)=z}A_w$. In particular, if indistinguishable worlds disagree on whether any implementation exists, no evidence-only procedure can always decide feasibility correctly and produce an implementation whenever one exists.
\end{proposition}
\begin{proof}
An evidence-only procedure has the same input, and hence the same chosen output, in every world of the fibre. Soundness requires validity in all of them. If feasibility differs within a fibre, either answer is wrong in one of its worlds.
\end{proof}

A native construction theorem therefore needs independent access and model premises. The positive results below supply explicit algorithms under such premises rather than replacing them by a requirement that the desired interface already be realizable.

\section{The general categorical mechanism and its sharp boundary}
Let $\mathcal C$ be a category of admitted native programs with an actual runtime interpretation $U:\mathcal C\to\mathbf{Set}$. The arrows of $\mathcal C$ are supplied executable operations, not all set maps merely asserted to exist. Program equality here means the equality appropriate to that declared semantic category; proof or execution histories may be retained separately. A morphism $h:A\to B$ is \emph{regular} in the present algebraic sense if there is a native $b:B\to A$ with $hbh=h$. This use of ``regular'' is not the regular-epimorphism terminology of regular categories.

\begin{theorem}[Native regular interactions generate retract roles]\label{m-thm:abstract}
From executable $h,b$ and the checked equation $hbh=h$, form $e=hb$. Then:
\begin{enumerate}
\item $e$ is idempotent, and the runtime set
\[
       P_h=\{y\in UB:Ue(y)=y\}=\im(Uh)
\]
is implemented by the wrapper $y\mapsto Ue(y)$. The normalized preparation operation $\bar b=bhb$ maps it back to $UA$ and satisfies $h\bar b=e$.
\item In $\Kar(\mathcal C)$, the new role $(B,e)$ gives a factorization of $h$ into a split epimorphism followed by a split monomorphism. Conversely, an arrow admitting such a factorization through an old object has a native regular inverse.
\item The role is isomorphic to an old object if and only if $e$ already splits in $\mathcal C$. If $\mathcal C$ is idempotent complete, this mechanism cannot create a new object up to the admitted isomorphism.
\item Native operations between generated roles are exactly arrows $f$ with $f=dfe$, and they compose without choosing representatives. If $\mathcal C$ is symmetric monoidal, tensor extends by $(A,e)\otimes(B,d)=(A\otimes B,e\otimes d)$ with inherited coherence.
\item Replacing a regular inverse $b$ by any other such inverse produces a canonically isomorphic generated role, with comparison arrows constructed from the two projectors.
\end{enumerate}
\end{theorem}
\begin{proof}
The equation $e^2=hbhb=hb=e$ is immediate. Every output $Uh(a)$ is fixed since $eh=h$, and every fixed $y$ equals $Uh(Ub(y))$, proving the set equality. The wrapper needs no membership test. The equation $h\bar b=hbhb=hb=e$ supplies a native preparation. Also $\bar b e=bhbhb=bhb=\bar b$, so it is typed on the generated role.

In the completion, $h:(A,1)\to(B,e)$ is split epic, with section $\bar b$. The inclusion $(B,e)\to(B,1)$ and its retraction are both $e$, so are split monic and split epic respectively. Their composite with $h$ is the original arrow. Conversely, if $h=ip$ with $ps=1$, $ri=1$, then $b=sr$ gives $hbh=ipsrip=ip=h$.

An isomorphism from $(B,e)$ to $(C,1)$ is exactly a pair of old arrows $i:C\to B$, $r:B\to C$ with $ir=e$, $ri=1$; this is precisely a splitting. The composition and tensor statements follow from $f=dfe$, as in the explicit matrix proof below, and the inherited associators and symmetries respect idempotents by naturality.

For a second inverse $b'$, put $e'=hb'$. Then $e'e=e$ and $ee'=e'$. The arrow represented by $e$ from $(B,e)$ to $(B,e')$ and the arrow represented by $e'$ in the reverse direction are inverse. These comparisons are independent of any pointwise choices and compose as the unique comparisons preserving the inclusions into $B$.
\end{proof}

The last uniqueness assertion concerns comparisons preserving the displayed inclusions: if $t:(B,e)\to(B,e')$ satisfies $e't=e$, then $t=e't=e$. This avoids claiming that there is only one abstract isomorphism between the roles.

This theorem is classical idempotent-splitting infrastructure with a runtime interpretation. It is not yet a method to \emph{find} $b$. The matrix construction below supplies that missing construction in a broad native class, from independently checkable response conditions. Requiring a native $b$ in every application without producing it would merely move the realization problem into a premise. Conversely, describing only a formal retract in a presheaf category would not establish its availability in $\mathcal C$ or in the runtime.



\subsection{Retained context is pullback to a resource base}
The second categorical construction concerns \emph{use after acquisition}, not acquisition itself. It applies independently of the retract construction.

\begin{theorem}[Retained-context and slice correspondence]\label{thm:context}
Let $\mathcal D$ have finite products and let $T\in\mathcal D$. Define $\mathcal D_T$ to have the same object labels as $\mathcal D$ and
\[
 \mathcal D_T(A,B)=\mathcal D(T\times A,B),\qquad
 g\star f=g\langle\pi_T,f\rangle,
 \qquad 1_A=\pi_A.
\]
This is a category. It is canonically isomorphic to the category of projections $T\times A\to T$ labelled by $A$, with all slice arrows between them, by
\[
 f\longmapsto\widehat f=\langle\pi_T,f\rangle.
\]
Forgetting repeated labels gives an equivalence with the actual full subcategory of $\mathcal D/T$ on those projections. A resource map $u:T'\to T$ induces a functor $u^*: \mathcal D_T\to\mathcal D_{T'}$ by precomposition with $u\times\id$; these functors compose strictly.

If $\mathcal D=\mathcal A/X$ for a category $\mathcal A$ with finite limits, products are fibre products over $X$. Thus contextual use is exactly a map over the token object $T\to X$ between the pullbacks $T\times_X A$ and $T\times_X B$.
\end{theorem}
\begin{proof}
An arrow $k:T\times A\to T\times B$ over $T$ has first component $\pi_T$, so is uniquely $\langle\pi_T,f\rangle$ for $f=\pi_Bk$. Composition of two such arrows is
\[
 \langle\pi_T,g\rangle\langle\pi_T,f\rangle
 =\langle\pi_T,g\langle\pi_T,f\rangle\rangle.
\]
Identity is $\langle\pi_T,\pi_A\rangle$. Associativity and the category laws follow from those in the slice, and the displayed inverse proves the correspondence. Substituting $u(t')$ in both stages proves the resource functor laws. Products in $\mathcal A/X$ are by definition the indicated pullbacks; also $(\mathcal A/X)/T$ is canonically $\mathcal A/T$, proving the final assertion.
\end{proof}

The token remains present; it is not independently resampled between composed arrows. This is the ordinary cartesian coKleisli construction. In a linear, quantum, destructive, or cost-sensitive runtime, cartesian sharing and exact retention require their own native justification. The theorem creates no token and cannot decide whether one is physically obtainable.

There are two precise instantiations below. For group actions, $\mathcal D=G\text{-}\Set$ and $T=G/H$; evaluation at its distinguished coset identifies the resulting arrows with $H$-equivariant maps. For geometry, $\mathcal D=\mathbf{Top}/X$ and $T=\Fr(P)\to X$; maps used for vector operations are restricted to be linear in the value fibres, although they need not be linear in the frame token. The frame bundle is not a vector bundle object with linear arrows, and the whole acquisition protocol is not an arrow of the fixed matrix category. These two instances share the proved pullback/context mechanism, not a theorem identifying their complete native languages.

The retract construction and the retained-context construction are likewise not asserted to commute in full generality. Context-dependent idempotents can encode additional objects. What is used here is explicit: first construct the role in the declared linear category, then pull its actual bundle and operations back to the acquired token base.

\section{Source access for the constructive mechanism}\label{sec:source}
Let $X$ be a nonempty compact Hausdorff space of public operating conditions. Let $R$ be a commutative unital subring of $C(X,\R)$ containing the rational constants. Each supplied element of $R$ is represented by a finite coefficient program with its own evaluation specification. The language admits finite addition and multiplication, transpose and finite matrix assembly, and reciprocal of a coefficient function accompanied by a proof that it is nowhere zero. The ring $R$ includes the resulting certified reciprocals; if one starts with a smaller coefficient ring, take the declared localization before fixing the old language. These operations generate further coefficient programs in $R$. This is an exact arithmetic model; finite-precision implementation requires a separate error analysis.

Old roles are the finite free modules $R^k$, represented geometrically by trivial parameterized vector spaces $X\times\R^k$ with operations restricted to $R$-coefficient matrices. Old operations are finite matrices over $R$, executed pointwise at the current $x\in X$. This defines the category $\Mat(R)$, with objects natural numbers and composition matrix multiplication. Sequential and parallel operations are supplied natively, not reconstructed from a category of models. Parameter $x$ is retained through a wire. An operation that changes $x$ requires additional transport data and is not silently included.

The input interaction is a supplied or identified matrix
\[
                         H\in\Mat_{n\times m}(R).
\]
Operationally, $u\in\R^m$ is a preparation or incoming signal and $H(x)u\in\R^n$ is the joint response at condition $x$. Linearity and the coefficient model need their own source warrant. Merely observing a few responses does not establish either assumption.


The distinction between a global coefficient program and pointwise access to an unknown device will remain explicit. Section~\ref{sec:identification} supplies a finite identification theorem under a finite-dimensional response promise; without that promise the point-query impossibility theorem applies. For frame acquisition, exact sampling, comparisons, retention, and repetition are additional native primitives, declared at the point they are used.

\section{Construct the role from the interaction, not from a chosen splitting}
\subsection{Finite rank certificates}
Set $G=HH^T\in\Mat_n(R)$. Write $c_j(G)$ for the $j$th elementary symmetric polynomial in its eigenvalues, defined algebraically as the sum of its principal $j$-minors; put $c_0=1$. A constant-rank-$r$ certificate, for $1\le r\le\min(m,n)$, consists of
\begin{equation}\label{m-eq:rank-certificate}
 \text{all $(r+1)$-minors of $H$ vanish},\qquad c_r(G)\ge\delta>0
 \quad\text{on }X.
\end{equation}
If there are no $(r+1)$-minors, the first condition is empty. By Cauchy--Binet,
\[
 c_r(HH^T)=\sum_{\substack{|I|=r\\|J|=r}}\det(H_{I,J})^2.
\]
Thus the certificate proves rank exactly $r$. Conversely, continuous constant rank on compact $X$ implies the existence of some $\delta>0$, but does not automatically provide an algorithm to find it. In a concrete symbolic or semialgebraic class, the equalities and positive lower bound can be independently checked. For example, polynomial entries with real algebraic coefficients on a compact semialgebraic $X$ admit these checks by real quantifier elimination. That decidability is not asserted for arbitrary coefficient programs.

Define the finite polynomial in $G$
\begin{equation}\label{m-eq:Q}
 Q=\frac{1}{c_r}\sum_{j=0}^{r-1}(-1)^j c_{r-1-j}G^j,
 \qquad e=GQ,\qquad B=H^TQ.
\end{equation}
For $r=0$, the certificate is $H=0$ and the definitions are $e=0$, $B=0$.

\begin{theorem}[Executable image generation]\label{m-thm:image}
From $H$ and \eqref{m-eq:rank-certificate}, the formulas \eqref{m-eq:Q} construct native coefficient programs satisfying
\[
 e^2=e=e^T,\qquad HB=e,\qquad eH=H,\qquad HBH=H,
 \qquad BHB=B.
\]
At every $x$, $e(x)$ is the orthogonal projector onto $\im H(x)$ and $B(x)$ is the Moore--Penrose inverse of $H(x)$. In particular,
\[
 P_H=eR^n=\im(H:R^m\to R^n)
\]
is constructed as a finitely generated projective module. The two maps
\[
 R^m\xrightarrow{\widehat H} P_H\xrightarrow{\iota}R^n
\]
implement the original interaction, with $\iota\widehat H=H$. The map $B:P_H\to R^m$ supplies a preparation for every admitted response, since $\widehat HB=\id_{P_H}$.

No eigenbasis, local frame, or solution of a pointwise optimization problem is an input to the algorithm.
\end{theorem}
\begin{proof}
At a fixed $x$, $G(x)$ is symmetric positive semidefinite with positive eigenvalues $\lambda_1,\ldots,\lambda_r$ and otherwise zero. The polynomial
\[
 p(t)=t^r-c_1t^{r-1}+c_2t^{r-2}-\cdots+(-1)^r c_r
\]
vanishes on each $\lambda_i$. Rearrangement gives
\[
 \lambda_i\big(c_{r-1}-c_{r-2}\lambda_i+\cdots+(-1)^{r-1}\lambda_i^{r-1}\big)=c_r.
\]
Therefore $GQ$ acts as identity on the positive eigenspace and as zero on the kernel. It is the asserted orthogonal projector. The identities $e^2=e=e^T$, $HB=e$, and $eH=H$ follow pointwise, hence as functions in $R$.

On $\im G$, $Q$ is inverse to $G$. Since $H^T$ kills $\ker G$, $H^TQ$ is inverse in the usual singular-vector decomposition on nonzero singular directions and zero on the orthogonal complement. Hence $BHB=B$ and the two symmetry conditions for the Moore--Penrose inverse hold. This argument proves the formula; the algorithm itself does not compute the decomposition.

The equation $eH=H$ gives $\im H\subseteq eR^n$. The reverse inclusion follows from $e=HB$, globally over $R$, not only fibrewise. Finally
\[
 R^n=eR^n\oplus(1-e)R^n
\]
by $v=ev+(1-e)v$ and the vanishing intersection. Thus $P_H$ is a finite projective module, and all displayed maps are restrictions of constructed matrices. Their identities are the asserted identities on the invariant domain.
\end{proof}

\begin{remark}[Cost and conditioning]
All formulas use finitely many ring operations and the one certified coefficient reciprocal $c_r^{-1}$. The $c_j$ can be computed either from minors or by Newton identities
\[
 c_k=\frac1k\sum_{i=1}^k(-1)^{i-1}c_{k-i}\tr(G^i).
\]
Using ordinary dense multiplication, one coarse arithmetic bound for forming $G,Q,e,B$ at a fixed parameter is $O(mn^2+rn^3+r^2)$; coefficient-evaluation and bit-complexity costs are additional. A stored valid value of $P_H$ uses $n$ scalar coordinates. Introduction by $e$ and use by a stored $a\times n$ matrix cost $O(n^2)$ and $O(an)$ scalar arithmetic operations respectively.

If $\|H(x)\|\le M$ and $c_r\ge\delta$, then the least nonzero singular value satisfies $\sigma_{\min}^2\ge\delta/M^{2(r-1)}$, so $\|B(x)\|\le M^{r-1}/\sqrt\delta$. The projector itself has norm $1$ when $r>0$. This does not make the polynomial formula numerically optimal; it exposes a nondegeneracy bound instead of hiding a division by a possibly vanishing determinant.
\end{remark}

\subsection{Why this is a native abstract type}
The new runtime role stores an $n$-vector and its type invariant $ev=v$. It has:
\begin{enumerate}
\item a constructor $\mathsf{pack}_e(v)=ev$;
\item an inclusion into the old representation $\mathsf{unpack}_e(v)=v$;
\item a constructor from actual interaction output, $u\mapsto Hu$;
\item a preparation operation $v\mapsto Bv$ for the generated response role.
\end{enumerate}
The invariant is established by construction; no runtime equality oracle is needed. An arbitrary invalid vector is not silently accepted as a value of the new type. The identity of this role is the matrix $e$, which acts as identity \emph{on valid values}; erasing a role to $R^n$ does not send that matrix to $I_n$. Consequently this erasure is not an identity-preserving functor to the old category, and calling it one would incorrectly erase the novelty.

\section{Coherent reuse and further role formation}
For an idempotent $e\in\Mat_n(R)$, write $(n,e)$ for its generated role. A native operation
\[
                         (n,e)\longrightarrow(m,d)
\]
is a matrix $F\in\Mat_{m\times n}(R)$ with $F=dFe$. Given any old matrix $A$, the sandwich $dAe$ automatically meets that requirement. Composition is multiplication and identity at $(n,e)$ is $e$.

\begin{theorem}[Compositional realization and higher-order roles]\label{m-thm:reuse}
The generated roles and these operations form $\Kar(\Mat(R))$. Interpretation by $eR^n$ is fully faithful into $R$-modules and has as its essential image the finite projective modules. It extends the old role category fully faithfully.

The following constructions have native matrix implementations and preserve the stated identities:
\begin{enumerate}
\item Parallel direct sum: $(n,e)\oplus(m,d)=(n+m,e\oplus d)$.
\item Tensor: $(n,e)\otimes(m,d)=(nm,e\otimes d)$.
\item The role of operations from $(n,e)$ to $(m,d)$: the image of the idempotent
\[
 \Lambda_{e,d}:\Mat_{m\times n}(R)\to\Mat_{m\times n}(R),\qquad A\mapsto dAe.
\]
It is represented, after a fixed column-vectorization convention, by $e^T\otimes d$. Evaluation is matrix--vector multiplication and composition is matrix multiplication. Their bilinearity gives the appropriate arrows out of tensor products.
\item A further idempotent $p:(n,e)\to(n,e)$ already gives the ordinary generated role $(n,p)$, since $p=epe$ and $p^2=p$. Hence iterating this role-formation mechanism adds no objects beyond the same completion.
\end{enumerate}
Sequential and parallel substitutions, evaluation, and composition obey their usual coherence laws. Every finite $R$-linear tensor network with only old boundary roles compiles to one old matrix; creation of its intermediate roles adds no new old-boundary linear function.
\end{theorem}
\begin{proof}
If $F=dFe$ and $G=fGd$, then $GF=fGdFe=fGFe$, so composition is well typed. Also $dF=F=Fe$, proving the identity laws; associativity is matrix associativity.

Restriction sends $F$ to an $R$-linear map $eR^n\to dR^m$. Conversely, any such map $\alpha$ extends uniquely to a sandwich matrix by
\[
 R^n\xrightarrow{e}eR^n\xrightarrow{\alpha}dR^m\hookrightarrow R^m.
\]
Restriction and this extension are inverse. A finite projective module is a direct summand of a finite free one, whose retraction gives the idempotent; the converse was proved above. On $(n,I_n)$ and $(m,I_m)$ the sandwich condition imposes no restriction, proving full faithfulness of the old inclusion.

Direct sums and Kronecker products preserve idempotence. The image of $e\otimes d$ is naturally $eR^n\otimes_R dR^m$, using the two displayed direct-summand decompositions. Further, $\Lambda_{e,d}^2=\Lambda_{e,d}$, and its fixed matrices are exactly the already identified hom-module. For a fixed matrix $F=dFe$ and valid $v=ev$, the equation $dFv=Fv$ proves typed evaluation. Products of fixed matrices are fixed at their outer endpoints, proving typed composition. All associativity, symmetry, unit, and evaluation identities are those of finite matrices and their tensor products; no choices of local frames enter them.

For an iterated idempotent $p$, its inclusion and retraction inside the completion are both represented by $p$, and their composite at the smaller role is its identity $p$. This is precisely the claimed collapse of repeated completion. Finally, replacing each internal operation by its actual sandwich matrix and each wire by ordinary contraction gives a finite old matrix expression at the boundary. Induction over the network proves equality of its interpretation. Resource costs are the costs of that expression, not a claim that internal abstract roles are free.
\end{proof}

The matrix-compilation clause concerns correctly typed linear and tensor wiring. Turning a pair of independently supplied vectors into their tensor product is bilinear, and copying $v$ into $v\otimes v$ is nonlinear; neither is silently a linear arrow out of the original vector role. A richer arithmetic runtime can execute such operations, but then its own program language, not $\Mat(R)$ alone, is the baseline for Theorem~\ref{thm:closure}.

This answers part of the request that a generated interface enable further interactions: it supports its own actual operations, operations \emph{as values}, typed evaluation, and subsequent image formation. What is new is the admissible intermediate type and its compositional specification. The theorem deliberately rules out the stronger claim that a type wrapper creates unimplementable old-type computations.

\subsection{Changing presentation without changing the generated role}
Suppose two interactions $H$ and $K$, in the same ambient output role, have the same image module. Their orthogonal projectors coincide pointwise and hence as functions. More generally, any supplied pair of sandwich matrices $U,V$ with $VU=e$ and $UV=d$ gives a native isomorphism between the generated roles, independently of the presentations used to find them. Such isomorphisms compose by matrix multiplication. This is coherent substitution, not a choice of representatives asserted to commute without proof.

\section{An exact old-language definability test}
\begin{proposition}[What an old codec would have to provide]\label{m-prop:codec}
The generated role $(n,e)$ can be represented by an old free role $R^k$, preserving its introduction, elimination, and all native uses, if and only if there are matrices
\[
 U\in\Mat_{n\times k}(R),\qquad V\in\Mat_{k\times n}(R)
\]
satisfying
\begin{equation}\label{m-eq:codec}
 eU=U,\quad Ve=V,\quad UV=e,\quad VU=I_k.
\end{equation}
Equivalently, $eR^n$ is a free module of rank $k$. Equivalently, the functor on old roles
\[
 F_e(R^a)=\{A\in\Mat_{n\times a}(R):eA=A\}
\]
is representable by an old role.
\end{proposition}
\begin{proof}
The first equations say the proposed encoder and decoder are typed; the last two say they are inverse. Restriction gives a module isomorphism, and any module isomorphism extends by $e$ to matrices with these equations. The functor $F_e$ is a retract of the representable functor $\Mat(R)(-,R^n)$. If represented by $R^k$, the natural isomorphism gives the same retract splitting by Yoneda, hence matrices $U,V$ with \eqref{m-eq:codec}. Conversely those matrices define the required natural isomorphism by postcomposition.
\end{proof}

A global obstruction to \eqref{m-eq:codec} therefore prevents a reusable old-language representation. It does not prevent the redundant $n$-coordinate implementation already constructed. Nor does it exclude nonlinear, discontinuous, parameter-forgetting, or multi-chart interfaces if those are admitted by a different old language. The declared preservation and access conditions determine the statement.

\subsection{A complete finite certificate for rank-one roles}
Assume $e=e^T$ has rank one at every point. Supply a finite open cover $U_1,\ldots,U_N$ with the following data:
\begin{enumerate}
\item every $U_i$ and every nonempty $U_i\cap U_j$ is connected;
\item a vector $a_i\in R^n$ with $ea_i=a_i$ that is nowhere zero on $U_i$;
\item coefficient functions $\phi_i\in R$, with $\phi_i\ge0$, $\sum_i\phi_i=1$, and $\supp\phi_i\subseteq U_i$;
\item the finite overlap graph, a sample point in each nonempty pairwise overlap, and exact signs of $a_i^Ta_j$ at those points.
\end{enumerate}
The finite cover and support assertions are certificates with their own verification method, not inferred from a plotted picture. In semialgebraic implementations they can be given explicitly and checked by the same decision machinery. Local generators can always be chosen from columns of $e$ on the charts where that column is nonzero; existence of an appropriate \emph{native} subordinate partition is a further hypothesis on the coefficient language.

For each overlap define
\[
             \sigma_{ij}=\operatorname{sgn}(a_i^Ta_j)\in\{+1,-1\}.
\]
These signs are constant on the connected overlap: two nonzero vectors in a one-dimensional real fibre have nonzero inner product, and its sign is continuous. On a triple overlap they obey $\sigma_{ij}\sigma_{jk}=\sigma_{ik}$.

\begin{theorem}[Finite obstruction or compiled scalar interface]\label{m-thm:sign}
Under these finite certificate hypotheses, exactly one of the following occurs:
\begin{enumerate}
\item There are signs $\epsilon_i\in\{\pm1\}$ with
\[
                       \sigma_{ij}=\epsilon_i\epsilon_j
\]
on all edges. They can be found by graph traversal. Then
\[
 v=\sum_i\phi_i\epsilon_i a_i,\qquad
 U=v,\qquad V=\frac{v^T}{v^Tv}
\]
are native coefficient programs satisfying \eqref{m-eq:codec} for $k=1$. Thus the generated role has an executable old scalar representation.
\item There is a cycle in the overlap graph with product of edge signs $-1$. The cycle, chart data, and sign evaluations form a finite certificate that no old free role represents $(n,e)$.
\end{enumerate}
The graph step takes $O(N+E)$ sign operations, where $E$ is the number of overlap edges. It decides this class completely once its chart and coefficient certificates are supplied.
\end{theorem}
\begin{proof}
A consistent propagation of $\epsilon_i$ from one vertex in each connected component either assigns all signs or encounters an edge contradicting the previous assignment. In the latter case the two traversal paths and that edge give a closed walk of negative product; deleting repeated segments yields a negative cycle. Conversely every sign assignment has positive product on every cycle.

When signs exist, all nonzero vectors $\epsilon_i a_i(x)$ with $\phi_i(x)>0$ point in the same direction in the line $e(x)\R^n$. They have a convex combination $v(x)$ with at least one positive coefficient, so $v(x)\ne0$. It follows that $v^Tv$ is strictly positive everywhere and, by compactness, bounded below. Its certified reciprocal is allowed. Since $e$ is the orthogonal rank-one projector with image spanned by $v$, it equals $vv^T/(v^Tv)$. The codec equations follow.

Conversely, any scalar codec provides a nowhere-zero global vector $v\in eR^n$. On each connected $U_i$, the sign comparing $v$ to $a_i$ is constant, giving signs $\epsilon_i$ satisfying every overlap equation. A negative cycle makes that impossible even among all continuous linear codecs, hence certainly among native ones. A codec to $R^k$ would have $k=1$ by fibrewise rank, so no other old free dimension can evade the obstruction.
\end{proof}

This is more than an assumption that a non-definability certificate exists: the certificate is extracted by a finite algorithm from locally checked interaction data. The geometry supplies the finite charts; the algebraic program supplies the operations; the sign calculation answers representability. These obligations are not interchangeable.

\section{One interaction through all stages}\label{sec:mobius}
Let $X=S^1=\{(c,s):c^2+s^2=1\}$ and take a coefficient language containing the coordinate functions $c,s$. Consider the native interaction
\begin{equation}\label{m-eq:mobius}
 H(c,s)=\frac12
 \begin{pmatrix}
  1+c&s\\ s&1-c
 \end{pmatrix}.
\end{equation}
Direct multiplication using $c^2+s^2=1$ proves $H^2=H=H^T$ and $\tr H=1$. Its rank certificate is exact: rank is one and $c_1(HH^T)=1$. The construction returns $e=H$ and $B=H$.

\begin{theorem}[Native realization without an old global scalar role]\label{m-thm:mobius}
The two-port interaction \eqref{m-eq:mobius} constructs a rank-one reusable role with native introduction, preparation, elimination, and all operations of Theorem~\ref{m-thm:reuse}. No global continuous linear codec to any old free role exists. This obstruction holds even if every continuous coefficient function were admitted, not only the smaller supplied coefficient grammar.
\end{theorem}
\begin{proof}
All positive claims are already constructed. For non-definability, write $(c,s)=(\cos\theta,\sin\theta)$ for $0\le\theta\le2\pi$ and
\[
 w(\theta)=(\cos(\theta/2),\sin(\theta/2))^T.
\]
Then $e(\theta)=w(\theta)w(\theta)^T$. A scalar codec would give a continuous nowhere-zero periodic vector $v(\theta)$ in the image line. It has the form $v(\theta)=a(\theta)w(\theta)$, where $a=w^Tv$ is continuous and nowhere zero. But $w(2\pi)=-w(0)$ and $v(2\pi)=v(0)$ imply $a(2\pi)=-a(0)$. The intermediate value theorem forces a zero, contradiction. Fibre rank excludes all $k\ne1$.
\end{proof}

This is the M\"obius line bundle, not a newly discovered topological object. The point is that the projector is an actual finite two-port interaction program, the new role is compiled from its response, and the failure of an old global coordinate is proved. Selecting an arbitrary value at each $x$ does not solve the problem, because an old operation must vary continuously with $x$.

\subsection{A finite sign certificate in the same example}
Use three arcs with angular centers $0,2\pi/3,4\pi/3$, each of radius $\alpha$ with $\pi/3<\alpha<2\pi/3$, chosen so no triple overlap occurs. To keep every pairwise overlap connected, for example take $\alpha=\pi/2$; the first arc wraps around $0$. Let
\[
 w_i=(\cos(\theta_i/2),\sin(\theta_i/2)),\qquad a_i=e w_i.
\]
For this displayed certificate, admit the indicated fixed real algebraic constants (in particular $\sqrt3/2$), or use the same vectors only as continuous witnesses for the topological obstruction. The interaction itself requires only rational constants and $c,s$. On the respective arc, $a_i$ is nonzero. The signs on the successive overlap edges $(1,2),(2,3),(3,1)$ are $+,+,-$. Indeed, near the overlap of the first two lifted arcs both half-angle differences are less than $\pi/2$ in magnitude, as for the second and third. On the wrap-around overlap, the continuous half-angle lift for the third chart differs by sign from that for the first chart. Equivalently, one may evaluate at angles $\pi/3,\pi,5\pi/3$ to obtain those signs directly. The cycle product is $-1$, giving Theorem~\ref{m-thm:sign}'s finite negative certificate. A subordinate partition is only needed for the positive codec branch, not for this negative certificate.

\subsection{A new use whose result is old}
The metric pairing
\[
                      P_H\otimes_R P_H\longrightarrow R,
                      \qquad v\otimes w\longmapsto v^Tw
\]
is an isomorphism for this line role. An explicit inverse sends $1$ to
\[
                         \sum_{j=1}^2(eu_j)\otimes(eu_j),
\]
where $u_j$ are the standard basis vectors. At each point choose a unit vector $w$ in the line only for the proof: $eu_j=w_jw$, so the sum is $w\otimes w$ and the pairing takes it to $1$. Both displayed programs use $e$ alone and are independent of this local sign choice. Thus two generated roles support a globally defined scalar interaction even though neither individual role has a global scalar coordinate. This is a concrete reusable interaction, not only a renamed response table.

\section{A subsequent experiment and a newly synthesized use}
The generated role does more than identify a vector bundle. It supplies a finite, globally valid preparation--test family for learning and using later interactions without first choosing a global basis.

For any projector $e\in\Mat_n(R)$, put $p_j=eu_j$ and let $\ell_j$ be the $j$th ambient coordinate restricted to $eR^n$. These are native, and
\begin{equation}\label{m-eq:dual-frame}
                  \sum_{j=1}^n p_j\ell_j=\id_{eR^n}.
\end{equation}
The $p_j$ need not be independent or even nonzero everywhere. Identity \eqref{m-eq:dual-frame}, not a hidden global frame assumption, is what makes them sufficient. For a role generated by $H$, each $p_j$ is actually preparable through the source interaction: choose input $Bp_j$ and obtain $HBp_j=p_j$.

\begin{theorem}[Learning later operations on the generated interface]\label{m-thm:later}
Let $T:eR^n\to dR^a$ be an unknown $R$-linear interaction that is actually callable on prepared generated-role values. At a fixed parameter, its $n$ responses $T(p_j)$ determine its full canonical sandwich matrix
\[
                    F=[T(p_1)\ \cdots\ T(p_n)]
\]
and every later response by $T(v)=Fv$. Globally, if each entry of $F$ belongs to a known $q$-dimensional response class with an accessible invertible evaluation matrix, $qn$ vector-response experiments identify a native program for $T$ exactly within that promise. The reconstructed program satisfies $dF=F=Fe$.

If $e$ has rank one and $T:eR^n\to eR^n$, the inferred scalar gain is $a=\tr(F)\in R$, with $F=ae$. If an independent certificate proves $|a|\ge\eta>0$, the program $a^{-1}e$ is a native inverse controller on the generated role. No global scalar coordinate for values of the role is required.
\end{theorem}
\begin{proof}
Equation \eqref{m-eq:dual-frame} gives $T(v)=\sum_jT(p_j)\ell_j(v)$, which is exactly multiplication by $F$. Since each output lies in $dR^a$, $dF=F$. Applying the same equation to $p_j$ gives $Fe=F$. Entrywise inversion of the supplied evaluation matrix gives the global coefficient programs, as in Proposition~\ref{m-prop:learn}.

On each one-dimensional fibre, an endomorphism is multiplication by one scalar. Its ambient sandwich extension is zero on $\ker e$, so its trace is that scalar because $\tr e=1$. Therefore $F=(\tr F)e$ pointwise, hence over $R$. The lower-bound certificate permits the reciprocal, and $(a^{-1}e)(ae)=e=(ae)(a^{-1}e)$ proves the inverse on the actual generated type.
\end{proof}

In the M\"obius example this procedure learns and inverts an unknown nonvanishing gain while the negative sign cycle still forbids replacing the role itself by a global scalar channel. The distinction is operational: a scalar \emph{endomorphism parameter} can exist without a scalar \emph{coordinate on all values}. The earlier two-copy metric pairing supplies another legal later interaction. By contrast, a continuous fibrewise-nonzero linear readout from the line role to an old scalar role would itself be an isomorphism and is ruled out by the same cycle certificate.

Nothing here grants an unknown physical call to $T$: its source access is explicitly supplied. The construction determines which preparations are sufficient, how to realize them using the earlier interaction, and how to compile the resulting learned inverse. If the original language already allowed all these old-coordinate experiments and arbitrary invariant tracking, the wrapper changes first-class typing and reuse, not the physical set of obtainable responses.



\section{Acquire an actual frame by executing the interaction}\label{sec:acquisition}
The preceding construction produces a role and its native representation. It does not yet produce a frame token. We now give a protocol that does, rather than assume that such a token arrives already built.

Assume $r>0$, and suppose independent certificates give
\[
 \sigma_{\min}^{+}(H(x))\ge\sigma>0,\qquad \|H(x)\|\le M
 \quad\text{for all }x.
\]
The rank certificate gives the first bound with $\sigma=\sqrt\delta/M^{r-1}$ when an operator-norm upper bound is also available. These constants can be supplied as rational conservative bounds; the algorithm does not require a square-root primitive merely because its analysis writes this formula.

At the current operating condition $x$, the native runtime can sample independent standard normal scalars, store real vectors, compute finite real arithmetic, compare scalars, execute $H(x)$ on an input vector, and repeat a failed trial. Parameter $x$ and accepted values remain unchanged during certified subsequent use. Exact Gaussian sampling and exact real arithmetic are the stipulated semantics, not a floating-point or bounded-time hardware claim.

Choose $\varepsilon>0$ and $L>\sqrt r\,\varepsilon/\sigma$. One proposal is an $m\times r$ standard Gaussian matrix $U$. First reject unless $\|U\|_F<L$; this uses the squared norm and no square root. Only then execute the $r$ column probes through $H(x)$ and retain the actual response matrix $V=H(x)U$. Accept exactly when
\begin{equation}\label{eq:accept}
                       V^TV-\varepsilon^2I_r>0.
\end{equation}
Positive definiteness is decidable in this exact arithmetic model by Sylvester's criterion, using the leading principal minors. On acceptance compute
\[
                  D_V=(V^TV)^{-1}V^T
\]
by the ordinary adjugate formula. All denominators are nonzero by the acceptance certificate.

\begin{theorem}[Bounded excitation and native frame acquisition]\label{thm:frames}
The protocol terminates almost surely. Uniformly in $x$, its number $N$ of Gaussian proposals satisfies
\[
 \Pr(N>j)\le(1-p_0)^j,
 \qquad \mathbb E N\le p_0^{-1},
\]
for the explicit positive constant below. It executes at most $rN$ column probes, each with input norm less than $L$. Its output is an actually observed frame $V$ for $P_x=\im H(x)$ and a native decoder satisfying
\[
 \im V=P_x,\quad D_VV=I_r,\quad VD_V=e_x,\quad
 \|D_V\|<\varepsilon^{-1},\quad
 \|V\|<ML,\quad \kappa(V)<ML/\varepsilon.
\]
Here $e_x$ is the orthogonal image projector and $\kappa(V)$ is the ratio of largest to smallest nonzero singular values.

For example choose constants
\[
 \varepsilon/\sigma<a<L/\sqrt r,
 \qquad 0<\rho<\min(a-\varepsilon/\sigma,L-a\sqrt r).
\]
With $d=mr$ and $\operatorname{vol}_d(B_\rho)$ the Euclidean ball volume, one may take
\begin{equation}\label{eq:pbound}
 p_0=(2\pi)^{-d/2}
       \exp\bigl(-(a\sqrt r+\rho)^2/2\bigr)
       \operatorname{vol}_d(B_\rho)>0.
\end{equation}
This conservative lower bound is part of the proof, not a claim of efficient practical acquisition.
\end{theorem}
\begin{proof}
At a fixed $x$, take a singular-value decomposition only for analysis. Rotate the source coordinates so that the first $r$ coordinates span the nonzero right singular subspace. The rotated $U$ is still standard Gaussian, with blocks $[Z;W]$, where $Z$ is $r\times r$. In orthonormal target coordinates,
\[
                  H(x)U=O\Lambda Z,
 \qquad s_{\min}(H(x)U)\ge\sigma s_{\min}(Z).
\]
Orthogonal source rotation preserves the Frobenius norm. The event
\[
 s_{\min}(Z)>\varepsilon/\sigma,
 \qquad \|[Z;W]\|_F<L
\]
therefore implies acceptance. The Frobenius ball of radius $\rho$ about $[aI_r;0]$ lies in that event: the perturbation in operator norm is less than $\rho$, so $s_{\min}(Z)>a-\rho>\varepsilon/\sigma$, while the total norm is less than $a\sqrt r+\rho<L$. Throughout this ball the standard Gaussian density in $d$ dimensions is at least the density factor in \eqref{eq:pbound}. Integrating gives probability at least $p_0$. Its value is independent of $x$.

Independent proposals at the retained $x$ now give the geometric tail bound. Summing $\Pr(N>j)$ for $j\ge0$ gives the expected-proposal bound and implies almost-sure termination. Proposals rejected by their norm require no physical probe; every remaining proposal uses at most $r$ calls, each input column having norm less than $L$.

On acceptance $s_{\min}(V)>\varepsilon$, so $V$ has rank $r$. Its image is contained in, and has the same dimension as, $P_x$, hence equals it. The full-column-rank inverse formula gives $D_VV=I_r$ and makes $VD_V$ the orthogonal projector onto $P_x$. Also $\|D_V\|=1/s_{\min}(V)<1/\varepsilon$ and $\|V\|\le\|H\|\|U\|<ML$, proving the conditioning bound. The runtime did not compute a singular-value decomposition; it executed the displayed arithmetic and certificate tests.
\end{proof}

The result is stronger than arbitrary pointwise choice: it is an actual protocol with a tail bound and bounded physical excitations. It is weaker than a total bounded-time deterministic frame selector. The acquisition resources include the supplied interaction, Gaussian source, exact comparisons, persistent memory, and permission to repeat. Calling the output a new resource does not hide those inputs.

\subsection{The full frame bundle and the accepted subset}
Define
\[
 \Fr(P)=\{(x,V):V:\R^r\xrightarrow{\cong}P_x\},\qquad
 \pi(x,V)=x.
\]
Right multiplication $V\mapsto Vg$ makes $\Fr(P)$ a principal $\mathrm{GL}_r(\R)$-bundle: on every neighbourhood where $r$ chosen projector columns are independent, those columns give a local frame, and every other frame has a unique invertible coordinate matrix. The accepted frames form a certified subset of this bundle. The singular-value threshold subset is not invariant under all of $\mathrm{GL}_r$, so it is not itself asserted to be a principal bundle.

These are general full frames, not necessarily orthonormal frames. An orthonormal/$O(r)$ formulation would require an additional native normalization algorithm, for example square roots and Gram--Schmidt. For a line, unnormalized nonzero vectors are not a single transitive $\{\pm1\}$-orbit: their magnitudes are invariant. The finite-group theorem below is not applied to that entire resource space by ignoring magnitude.

\begin{theorem}[Frame-dependent use and exact descent]\label{thm:frame-reuse}
The acquired token gives inverse codecs
\[
 \mathsf{encode}_V(z)=Vz,\qquad \mathsf{decode}_V(v)=D_Vv
\]
between $\R^r$ and the fibre $P_x$, retaining $(x,V)$. Globally these are a trivialization of the pullback $\pi^*P$ over $\Fr(P)$.

For native linear role maps $F:P_x\to Q_x$ and $K:Q_x\to R_x$, with retained frames $V,W,Z_0$ of the three roles, define
\[
 A_F=D_WFV,\qquad A_K=D_{Z_0}KW.
\]
Then $A_{\id}=I$ when an identity uses the same source and target frame, $A_KA_F=A_{KF}$, and direct-sum and tensor coordinates respect their corresponding compositions. Under $V'=Vg$ and $W'=Wh$,
\[
                         A'_F=h^{-1}A_Fg.
\]
Thus coordinate operations compose exactly when the same intermediate frame is retained.

For a continuous family of fibrewise-linear endomorphisms encoded by matrices $A(V)$ on the full frame bundle, a continuous old-base operation $F:P\to P$ exists with $A(V)=D_VF_xV$ if and only if
\begin{equation}\label{eq:descent}
                    A(Vg)=g^{-1}A(V)g
\end{equation}
for every frame and every $g\in\mathrm{GL}_r$. This is a continuous descent criterion, not an automatic native-program synthesis theorem.
\end{theorem}
\begin{proof}
The identities $D_VV=I_r$ and $VD_V=e_x$ prove that the codecs are inverse on valid values. The formulas are continuous on full frames, so they trivialize the pullback. For sequential reuse,
\[
 (D_{Z_0}KW)(D_WFV)
 =D_{Z_0}K(WD_W)FV=D_{Z_0}KFV,
\]
since $WD_W$ is identity on $Q_x$. The identity law is $D_VV=I$. For direct sums use block-diagonal frames and decoders; for tensors use $V\otimes W$ and $D_V\otimes D_W$. Matrix multiplication and the mixed-product rule prove the stated composition laws. Finally $D_{Wh}=h^{-1}D_W$, directly from the full-rank inverse formula, proving the change-of-frame law.

Necessity of \eqref{eq:descent} follows from that law. For sufficiency define, at any frame over $x$,
\[
                         F_x=V A(V)D_V\big|_{P_x}.
\]
Changing $V$ to $Vg$ leaves this expression unchanged by \eqref{eq:descent}. Hence it is well defined. Local continuous sections of the frame bundle make it continuous on each trivializing neighbourhood, and independence of frame makes those local definitions agree. This produces the continuous bundle map. No claim that its coefficient functions belong to the supplied native grammar follows without further evidence or a compiler.
\end{proof}

This is the geometric instance of Theorem~\ref{thm:context} in $\mathbf{Top}/X$, restricted to fibrewise-linear uses. Its descent test asks when retained context can be forgotten. The scalar decoder of a nontrivial role does not descend to an old global scalar codec; it depends essentially on the acquired frame. A continuous section of $\Fr(P)\to X$ exists exactly when $P$ has a continuous global frame, equivalently is a trivial bundle: a section is such a frame, and a trivialization supplies one.

\subsection{Complete execution of the circle example}
For the interaction in \eqref{m-eq:mobius}, $r=1$ and $\sigma=M=1$. Choose, for instance, $\varepsilon=1/4$ and $L=2$. Draw $u\in\R^2$ standard Gaussian, reject before probing if $\|u\|\ge2$, execute $v=H(x)u$, and accept if $v^Tv>1/16$. The output $v$ and decoder $v^T/(v^Tv)$ give an actual scalar interface for the image line at the retained $x$, with decoder norm below $4$. Theorem~\ref{thm:frames} gives a uniform finite expectation; the protocol performs no global sign selection.

The same role still has the negative sign cycle and no global continuous scalar codec on $S^1$. There is no contradiction: the acquired interface is over $(x,v)$, not over $x$ alone. A draw may fail at some other operating point, and the almost-sure assertion is for each retained $x$, not one draw simultaneously valid for all $x$. Changing $x$ requires reacquisition or independently justified transport. The protocol also does not make the original fixed linear language spontaneously contain sampling or looping.

The later gain experiment of Theorem~\ref{m-thm:later} now has two usable implementations. It can use the global redundant dual frame, or use the acquired scalar coordinates at a fixed $x$. Both compile to declared arithmetic and physical calls. The first demonstrates that first-class role formation does not require a global basis; the second demonstrates a physically acquired reference whose exact continued retention warrants simpler coordinates. Neither creates an unknown reverse-device oracle.

\section{What finite interaction evidence can establish}\label{sec:identification}
\begin{proposition}[Exact finite response acquisition under a finite model promise]\label{m-prop:learn}
Suppose every entry of $H$ is known a priori to lie in the span of known coefficient functions $\psi_1,\ldots,\psi_q$, and accessible parameter points $x_1,\ldots,x_q$ make
\[
                  E=(\psi_j(x_i))_{i,j}
\]
invertible with an exact native inverse. Suppose each standard input vector of $R^m$ can be prepared and the full output vector read exactly at each $x_i$. Then $qm$ vector-response experiments, or $qmn$ scalar readouts, identify all coefficient programs of $H$ by multiplication with $E^{-1}$. Given a checkable rank certificate in that coefficient class, Theorem~\ref{m-thm:image} then compiles its reusable role and operations.
\end{proposition}
\begin{proof}
At $x_i$, the response to the $j$th standard input is the $j$th column of $H(x_i)$. For each matrix entry the $q$ observed values are $E$ times its unknown coefficient vector. The inverse computes that vector. The supplied finite-dimensional promise makes the reconstruction globally correct; no extrapolation claim is inferred merely from fitting the samples. Rank verification and role compilation are subsequent finite steps under their separately stated hypotheses.
\end{proof}

For \eqref{m-eq:mobius}, the entry class is $\operatorname{span}(1,c,s)$. Parameters $(1,0),(0,1),(-1,0)$ yield an invertible evaluation matrix. Six two-coordinate responses identify the interaction within that declared class. The polynomial identity $c^2+s^2=1$ verifies the rank certificate and the displayed half-angle proof or finite sign cycle verifies non-definability. The observations alone, without the entry-class promise, do not force the result.

\begin{theorem}[Finite samples cannot recognize global role novelty]\label{m-thm:no-samples}
On the unrestricted class of continuous rank-one orthogonal projectors $S^1\to\Mat_2(\R)$, no deterministic always-terminating finite point-query procedure can decide whether the generated role has a global continuous scalar codec. This allows adaptive queries and exact projector values.
\end{theorem}
\begin{proof}
Run a putative procedure on the constant projector $e_0=\operatorname{diag}(1,0)$. It halts after querying finitely many points. Choose an open arc disjoint from those points. Construct a continuous angle $\beta:[0,2\pi]\to\R$ that is constant $0$ up to that arc, rises continuously from $0$ to $2\pi$ inside it, and remains $2\pi$ afterward, choosing the angular cut outside the arc. Set
\[
 \widetilde e(\theta)=\frac12
 \begin{pmatrix}1+\cos\beta(\theta)&\sin\beta(\theta)\\
 \sin\beta(\theta)&1-\cos\beta(\theta)\end{pmatrix}.
\]
This is a continuous periodic rank-one projector equal to $e_0$ at every queried point. The same half-angle proof as above shows its line role has no global scalar codec. The constant projector does have one. The adaptive transcript and hence the decision are identical, contradiction. The construction can be made smooth using a smooth transition on the unused arc.
\end{proof}

The impossibility concerns point queries without a global response-class or regularity-bound promise. It does not negate exact coefficient identification, interval certificates, a supplied transition cover, or a validated bound sufficient to exclude an unseen twist. It also does not say that topology is unstable under uniformly small errors: a uniform bound is much stronger evidence than finitely many unbounded point samples.

\section{Stability and exact limits of the formation mechanism}
\subsection{A native transport under uniformly small revision}
\begin{theorem}[Constructive stability of a generated role]\label{m-thm:stable}
Let $e,f\in\Mat_n(R)$ be symmetric idempotents with a verified uniform bound
\[
                           \|e(x)-f(x)\|\le\rho<1.
\]
Assume the coefficient language admits the certified reciprocal used in the adjugate formula below. Define
\[
 A=efe+(1-e),\qquad U=fe,\qquad V=A^{-1}ef,
 \qquad A^{-1}=\operatorname{adj}(A)/\det(A).
\]
Then $A$ is everywhere invertible, and $U:(n,e)\to(n,f)$ and $V:(n,f)\to(n,e)$ are native inverse isomorphisms. Thus this kind of sufficiently controlled evidence revision changes the projector presentation without changing the reusable role up to an explicit native isomorphism.
\end{theorem}
\begin{proof}
For $v\in\im e$, $\|fv\|\ge(1-\rho)\|v\|$, so $f|_{\im e}$ is injective. The symmetric argument makes $e|_{\im f}$ injective, hence finite fibre dimensions agree. Moreover
\[
 v^Tefe v=\|fv\|^2\ge(1-\rho)^2\|v\|^2.
\]
On $\ker e$, $A$ is identity. It is therefore positive definite, with eigenvalues bounded below by $(1-\rho)^2$, and $\det A\ge(1-\rho)^{2n}>0$. Its reciprocal has a global certificate. Since $A$ commutes with $e$, so does its inverse. Hence $V$ is typed and
\[
 VU=A^{-1}ef fe=A^{-1}efe=e.
\]
The map $U$ is a fibrewise injection between equal finite dimensions, hence a bijection. The last equation gives its inverse, so $UV=f$ pointwise, and therefore over $R$.
\end{proof}

These canonical pairwise maps are not claimed to compose strictly across a sequence of revisions: the chosen direct comparison $ge$ need not equal $(gf)(fe)$. Actual isomorphisms do compose coherently by multiplication, and distinct choices are related by automorphisms. This is the precise groupoid-level stability conclusion, rather than an unproved strictly functorial choice of projector transports.

\subsection{Rank jumps are a mathematical obstruction}
\begin{proposition}[No continuous image role across a rank jump]\label{m-prop:jump}
If $X$ is connected and $H(x)$ has nonconstant rank, no continuous idempotent matrix $e(x)$ of any fixed ambient size can have image isomorphic, fibrewise and continuously over $X$, to $\im H(x)$ at every point. In particular the exact generated-image role cannot be a finite vector bundle over all of $X$.
\end{proposition}
\begin{proof}
For any idempotent matrix, eigenvalues are $0$ and $1$, so its rank equals its trace. The trace is continuous and integer-valued, hence constant on connected $X$. A fibrewise image isomorphism would force it to equal the varying rank of $H$.
\end{proof}

For example $H(t)=[t]$ on $[-1,1]$ has image $\R$ away from $0$ and $0$ at $0$. The pointwise projector is $1$ off $0$ and $0$ at $0$, hence not continuous. One may use a stratified interface, a module that is not finite projective, or a deliberately approximate rank threshold. Each changes the object doctrine or approximation question; none is provided by this theorem. The obstruction shows why simply declaring the image a native role would be unsound.


\section{Exact capability change in the symmetry instance}\label{sec:symmetry}
Let $G$ be a finite group. A finite $G$-set models a raw interface whose presentations may be changed by $G$. Ordinary output symbols such as Boolean answers carry the trivial action. The baseline semantic operations are equivariant functions. Native implementation will be supplied separately in Section~\ref{g-sec:compiler}.

For a finite resource $G$-set $R$, define the category $\C_R$ by
\[
 \operatorname{Ob}(\C_R)=\operatorname{Ob}(G\text{-}\FinSet),\qquad
 \C_R(X,Y)=\Hom_G(R\times X,Y).
\]
The identity is $(r,x)\mapsto x$, and composition is
\[
 (q\mathbin{\star}p)(r,x)=q(r,p(r,x)).
\]
The operational implementation retains the resource:
\[
 \widehat p:R\times X\longrightarrow R\times Y,
 \qquad (r,x)\longmapsto(r,p(r,x)).
\]
This is Theorem~\ref{thm:context} for $G$-sets, equivalently the coKleisli category of the cartesian comonad $R\times-$. The cartesian structure shares one retained reference; it does not posit independently resampled references. It is a special context-dependent composition mechanism, not a replacement for the programme's richer contextad and native witness structure.

\begin{theorem}[Reference classification]\label{g-thm:reference}
For a subgroup $H\leq G$ and finite $G$-sets $X,Y$, there is a bijection
\begin{equation}\label{g-eq:classification}
 \Hom_G((G/H)\times X,Y)\cong\Hom_H(X,Y).
\end{equation}
Here $X,Y$ on the right have their restricted $H$-actions. The bijection is natural in equivariant maps of $X,Y$, preserves identity, sequential composition, products and typed substitution, and is explicitly computable from finite tables. It sends $p$ to $f(x)=p(H,x)$; its inverse is
\begin{equation}\label{g-eq:processor}
 p_f(gH,x)=g\,f(g^{-1}x).
\end{equation}
Consequently $\C_{G/H}$ is isomorphic to the category with the same object labels as $G\text{-}\FinSet$ and all $H$-equivariant functions between those objects. This does not assert that every $H$-set is the restriction of a $G$-set.
\end{theorem}
\begin{proof}
For $h\in H$, equivariance gives
$p(H,hx)=p(hH,hx)=h p(H,x)$, so evaluation at $H$ is $H$-equivariant. If $gH=ghH$, then
\[
 gh f(h^{-1}g^{-1}x)=g f(g^{-1}x),
\]
which proves that \eqref{g-eq:processor} is well-defined. It is $G$-equivariant because
$p_f(agH,ax)=ag f(g^{-1}x)=a p_f(gH,x)$.
The two formulas are inverse. For $f:X\to Y$ and $q:Y\to Z$,
\[
 p_q(gH,p_f(gH,x))
 =g q(g^{-1}g f(g^{-1}x))
 =g(qf)(g^{-1}x)=p_{qf}(gH,x).
\]
The identity case and the coordinatewise product and substitution cases use the same cancellation. Finite orbit and group tables make every displayed operation effective.\end{proof}

\begin{corollary}[Exact token elimination in the group instance]\label{g-cor:elimination}
For $f:X\to Y$ that is $H$-equivariant, the associated retained processor $p_f$ is independent of its token if and only if $f$ is $G$-equivariant. Thus the old globally equivariant operations are exactly the contextual operations that descend after forgetting the reference.
\end{corollary}
\begin{proof}
If $f$ is $G$-equivariant, then $g f(g^{-1}x)=f(x)$, so the formula for $p_f$ is independent of $gH$. Conversely suppose independence holds. For $a\in G$,
\[
 f(ax)=p_f(H,ax)=p_f(aH,ax)=a\,p_f(H,x)=a f(x).
\]
The middle equivariance is that of $p_f$, already proved.
\end{proof}

\begin{corollary}[A finite separating experiment]\label{g-cor:strict}
If $H'<H\leq G$, then the $H'$-equivariant repertoire on the inherited $G$-set objects strictly contains the $H$-equivariant repertoire. A witness is the experiment on the regular left $G$-set
\[
 t_{H'}:G\longrightarrow\Bool,
 \qquad t_{H'}(x)=[x\in H'].
\]
It is $H'$-invariant and not $H$-invariant. Thus any native old grammar whose generators are $H$-equivariant cannot implement this experiment without an additional resource or a changed observational contract.
\end{corollary}
\begin{proof}
Left multiplication by $H'$ preserves $H'$. If $h\in H\setminus H'$, then $t_{H'}(1)=1$ whereas $t_{H'}(h)=0$. Equivariance is preserved by the baseline structural operations.\end{proof}

This is a nondefinability argument against a closure property, not a failure to find a short old expression. It is also conditional on using the appropriate raw interfaces: if all interfaces carried trivial actions, no such symmetry obstruction would result.

\begin{proposition}[Coherent resource refinement]\label{g-prop:refinement}
An equivariant resource map $u:R'\to R$ induces a functor
\[
 u^*: \C_R\longrightarrow\C_{R'},\qquad
 p\longmapsto p\circ(u\times\id).
\]
For $R'=G/H'$, $R=G/H$ and $H'\leq H$, the quotient map induces precisely inclusion of $H$-equivariant operations among $H'$-equivariant operations under Theorem~\ref{g-thm:reference}. These functors compose strictly. A change of reference presentation by $g\in G$ transports a map by
\[
 T_g(f)(x)=g f(g^{-1}x),
\]
from the $H$-equivariant to the $gHg^{-1}$-equivariant repertoire, with $T_gT_h=T_{gh}$.
\end{proposition}
\begin{proof}
Substitution of $u(r')$ in the composition formula proves functoriality. Evaluation at the distinguished coset gives the subgroup assertion. The conjugation statement follows directly from equivariance and multiplication in $G$.\end{proof}

Thus different frames have coherent comparison maps, not a canonical preferred frame. If witness categories, costs or occurrence data are retained, the same formula at the level of functions does not produce the missing witness transformations. Those require a native lift.

\section{Interaction assembles a reference rather than naming one}
Suppose several actual interactions return tokens $r_i\in R_i$. Retain the tuple $r=(r_1,\ldots,r_m)$ rather than forgetting the earlier responses. Its orbit is a transitive resource, and
\begin{equation}\label{g-eq:intersection}
 \Stab(r)=\bigcap_{i=1}^m\Stab(r_i),
 \qquad G\cdot r\cong G/\Stab(r).
\end{equation}
The isomorphism is $g\Stab(r)\mapsto gr$. It is well-defined, bijective and equivariant by the orbit--stabilizer theorem. No assumption of independent stochastic responses is used.

\begin{theorem}[Finite reference formation]\label{g-thm:formation}
Let $R$ be a finite $G$-set with action kernel $K$. If an interaction protocol can actually acquire and retain tokens forming a tuple $r$ whose stabilizer is $K$, then that tuple constructs a reusable reference interface whose operations are exactly the $K$-equivariant operations on the old object labels. If $K=1$, every finite underlying function between those labels is enabled semantically. Such a tuple exists among the elements of $R$: the tuple of all elements has stabilizer $K$. A tuple with trivial stabilizer is a base of the action. If native finite table realization holds, the construction and its subsequent operations are executable. Existence of a base does not prove that the physical protocol can acquire it.
\end{theorem}
\begin{proof}
Intersecting the stabilizers of all elements gives precisely the kernel. Apply \eqref{g-eq:intersection} and Theorem~\ref{g-thm:reference}. The finite orbit identification and the processor tables are computable. The final clause uses the separately assumed native table interpretation.\end{proof}

A sharper finite selection bound is available when all candidate tokens have already been physically acquired and their finite action presentation is known. Start with $H=G$ and retain a candidate that strictly decreases $H$ whenever $H\ne K$. There is such a candidate, since the total intersection is $K$. Every strict decrease has index at least two, so at most $\lfloor\log_2[G:K]\rfloor$ retained candidates suffice. This is a bound on selecting a subtuple from an acquired finite collection, not a bound on the number of experiments needed to obtain that collection.

For $G=S_n$ acting on $n$ points, a tuple of $n-1$ distinct points is a base and no shorter tuple is. Its pointwise stabilizer permutes the unmarked points freely, so it is trivial exactly when at most one point is unmarked. This is a family of growing acquisition requirements, not one isolated example. However, a name-blind register machine with only equality, copying and selection cannot manufacture the one missing raw point from those $n-1$ registers. Thus a free orbit and the semantic formula \eqref{g-eq:processor} do not by themselves prove native realizability. The next section avoids that gap by giving a small explicit instruction set and compiler.

\begin{proposition}[No canonical resource creation]\label{g-prop:no-reference}
If $G\ne1$, there is no equivariant map $1\to G$ for the regular left action. More generally, an equivariant map cannot send a token with stabilizer $H$ to a token with a strictly smaller stabilizer. An equivariant deterministic acquisition process from invariant inputs therefore cannot create a complete reference from nothing.
\end{proposition}
\begin{proof}
If $h$ fixes an input $r$ and $u$ is equivariant, then $h u(r)=u(hr)=u(r)$, so $\Stab(r)\subseteq\Stab(u(r))$. Apply this to the invariant singleton input.\end{proof}

The environment or the retained joint input must carry the symmetry-breaking resource. Joint acquisition can reduce an investigator's residual stabilizer because the additional token was not in the investigator's former resource state. It does not violate the proposition for the complete environment-plus-investigator system. A random invariant sample, if retained and used relative to itself, can serve as a relative reference. It does not provide a reproducible canonical absolute orientation.

\section{A native executable family: independently calibrated polarities}\label{g-sec:compiler}
Fix $k\geq1$, let $G=(\mathbb Z/2)^k$, and let $X=\{0,1\}^k$ with the translation action. Coordinate $i$ has a raw bit sort $B_i$ on which only the $i$th factor flips polarity. Raw polarity values are not publicly named Boolean values. Public Booleans have the trivial $G$-action.

The baseline native language has:
\begin{itemize}
\item raw input registers, copying, discarding and finite products;
\item the toggle $b\mapsto b\oplus1$ on each $B_i$;
\item equality of two registers of the same raw sort, with a public Boolean answer;
\item public finite Boolean computation and branching;
\item finite stored public tables and finite program syntax.
\end{itemize}
It has no initialized raw register and no nullary raw constant. Every primitive is equivariant, as are every resulting program and its finite composition. In particular, a baseline experiment $X\to\Bool$ is constant because the action on $X$ is transitive. Every baseline unary $X\to X$ function, if defined, is a translation $x\mapsto x\oplus c$: equivariance gives $f(x)=x\oplus f(0)$.

\subsection{The acquisition protocol and physical accounting}
For each coordinate there is an external sample port that transfers one actual token $r_i:B_i$. The protocol opens each port once, stores the returned raw value in a persistent register, and closes the port. It then assembles $r=(r_1,\ldots,r_k)\in X$. Every possible resulting tuple has trivial stabilizer, so the procedure does not need to select a rare response or identify an absolute polarity.

The token is \emph{physically supplied}; persistent storage and the operations on raw registers are supplied capabilities. What is \emph{constructed} is a single reusable coordinate interface and a library of compiled operations relative to those retained tokens. A globally equivariant model represents a source port as carrying its raw token in its source state, and represents transfer by the identity from that resource wire into the receiver. It does not falsely model transfer as an equivariant map $1\to B_i$. The resources are relocated or copied as the stated source protocol permits; they are not created from an invariant singleton.

There is an explicit type boundary between acquiring a source sample and using a retained local reference. The baseline claim concerns clients with no reference registers. If a prior capability definition already lets a client acquire these samples and retain them for the whole experiment, Theorem~\ref{thm:closure} applies: there is then no strict increase over that larger old repertoire.

\begin{theorem}[Elementary relative-coordinate compiler]\label{g-thm:compiler}
After the $k$ transfers, the native language above implements, relative to its unchanged stored reference $r$, every function
\[
 f:X^n\longrightarrow X^m
\]
specified by a finite public table, as well as every finite public-valued experiment on $X^n$. The realized raw operation is
\begin{equation}\label{g-eq:relative-compiler}
 P_f(r;x_1,\ldots,x_n)
 =\bigl(r\oplus f_j(x_1\oplus r,\ldots,x_n\oplus r)\bigr)_{j=1}^m.
\end{equation}
Here $f_j$ is the $j$th $X$-valued output component. The compiler is effective, retains $r$, and preserves identity, sequential composition, products and typed substitution exactly at the level of functions. Its elementary finite instruction implementation proves native realizability, rather than assuming fullness of a semantic interpretation.
\end{theorem}
\begin{proof}
For each input coordinate, negating equality with the corresponding reference register computes the public relative bit $x_{j,i}\oplus r_i$. Public computation uses these $nk$ bits to look up the finite output table. For each output bit $b_i$, copy $r_i$ if $b_i=0$ and toggle a copy of $r_i$ if $b_i=1$. The reference registers are unchanged. These instructions yield \eqref{g-eq:relative-compiler}. Writing $\beta_r(x)=x\oplus r$, the map is $\beta_r^{-1}f\beta_r$ on each typed power. Therefore
$P_q(r;P_f(r;x))=P_{qf}(r;x)$, and products and substitution satisfy the same identity.\end{proof}

The table may be exponentially large in the arity. A direct implementation uses $nk$ raw equality tests, the chosen public table lookup, and $mk$ output copies or toggles. This is an existence and correctness compiler, not a promise of succinct descriptions, cheap hardware, statistical robustness or energetic advantage.

For example,
\[
 t_a(r,x)=[x=r\oplus a]
\]
is a native nonconstant experiment that was impossible as a baseline $X\to\Bool$ program. The family $(t_a)_{a\in X}$ separates all raw states relative to the retained reference. These are new available experiments, not merely new names for the old constant tests. The same compiler can prepare every raw state $r\oplus a$ from no live raw input. The resource is crucial both for the separating experiments and for these preparations.

Acquiring only coordinates in $J\subseteq\{1,\ldots,k\}$ leaves residual symmetry
\[
 H_J=\{g\in G:g_i=0\text{ for every }i\in J\}.
\]
Each new coordinate token strictly decreases this subgroup and enables an additional nonconstant relative-coordinate test. Thus the complete interface is built through a sequence of strictly growing resource-indexed repertoires. A token for one coordinate does not secretly resolve the other polarities.

\section{Formation after coupling versus separate interfaces}\label{g-sec:coupling}
The same mechanism gives an exact comparison with the earlier contextual-reduction problem. With a reference whose stabilizer is $H$, the available Boolean experiments on a finite $G$-set $X$ are precisely the $H$-invariant functions. Therefore two raw states give the same answer to every such experiment exactly when they belong to the same $H$-orbit. The resulting observational interface is the orbit set $H\backslash X$.

For two references with stabilizers $H,K$, the retained joint reference has stabilizer $H\cap K$. There is a canonical comparison
\begin{equation}\label{g-eq:coupling}
 c:(H\cap K)\backslash X\longrightarrow
 (H\backslash X)\times(K\backslash X),
 \qquad [x]_{H\cap K}\longmapsto([x]_H,[x]_K).
\end{equation}
Its domain records the experiments enabled by the joint resource. Its target records independent values in the separately formed observational interfaces. These are not automatically the same interface.

\begin{theorem}[Exact finite reference-coupling test]\label{g-thm:coupling}
The comparison \eqref{g-eq:coupling} is injective exactly when
\[
 Hx\cap Kx=(H\cap K)x\quad\text{for every }x\in X,
\]
and is surjective exactly when every $H$-orbit meets every $K$-orbit. Both conditions are decidable by finite action tables, without assuming a comparison isomorphism. For the left regular action on $X=G$, injectivity always holds and the comparison is bijective exactly when $HK=G$.
\end{theorem}
\begin{proof}
The raw states with the same pair of orbit labels as $x$ form $Hx\cap Kx$. The joint-resource label of $x$ is $(H\cap K)x$. Equality of those sets for every $x$ is precisely injectivity. A pair of orbit labels lies in the image precisely when the corresponding orbits have a common state, which proves the surjectivity assertion.

For the regular action, $Hx\cap Kx=(H\cap K)x$ by right cancellation. The domain has cardinality $|G|/|H\cap K|$ and the target has cardinality $|G|^2/(|H||K|)$. Since the map is already injective, it is bijective exactly when these cardinalities agree. The finite-group identity $|HK|=|H||K|/|H\cap K|$ gives the stated condition.\end{proof}

There can be synergy that separate quotient interfaces erase. Take $G=S_4$, $X=\{1,2,3,4\}$,
\[
 H=\langle(1234)\rangle,\qquad K=\langle(1243)\rangle.
\]
Both subgroups act transitively, so each separate orbit interface is a singleton. Their intersection is trivial: their nonidentity squares are different, and their order-four generators and inverses do not agree. Thus the jointly enabled interface has four states and \eqref{g-eq:coupling} is the map from four states to one. No independent pair of old orbit labels can recover the new jointly available tests.

Conversely, even when injectivity holds, independently chosen separate labels can be incompatible. For the regular action with $H=K<G$, the comparison is the diagonal of $H\backslash G$, so it is not surjective when that orbit set has more than one element. This is a physical shared-input constraint, not an implementation failure.

For the polarity family $G=(\mathbb Z/2)^k$ acting regularly, disjoint sets of acquired coordinate references give residual subgroups $H,K$ with $H+K=G$; the comparison is bijective. Overlapping acquired coordinate sets need not satisfy that condition. The finite subgroup calculation therefore supplies a concrete positive condition and an explicit obstruction for modularly composing the newly formed interfaces.

These are statements about the actual available operation families, hence require the appropriate native processor realization before being called physical capability theorems. The elementary polarity compiler discharges that condition. The general $S_4$ example is a semantic/native-table-model example unless an implementation for its raw carriers is supplied. Neither example asserts that an arbitrary native contextad has this quotient behavior.



\section{The combined result and what remains unresolved}
\begin{theorem}[A constructive staged formation mechanism]\label{thm:combined}
In the declared constant-rank native class, the following conclusions hold.
\begin{enumerate}
\item A response program with its rank certificate yields a first-class projective image role, a preparation right inverse, and native operations and operation-valued roles with coherent typed composition.
\item In the certified rank-one chart class, a finite algorithm either constructs an old scalar codec or returns a negative cycle excluding one. The circle interaction gives the negative case.
\item With the separately declared sampling, comparison, retention, and repetition primitives, bounded excitation of the source interaction acquires a frame almost surely with uniform expected cost and a bounded decoder. The resulting token is actual observed output, not an assumed choice of splitting or basis.
\item Continued use at the retained parameter is exactly frame-context use. Coordinates compose and change presentation by the proved gauge law. Complete finite-dimensional response promises support subsequent exact identification and the stated inverse-controller synthesis.
\end{enumerate}
In the separate finite symmetry implementation, actual acquisition and retention of polarity samples constructs a native reference interface. Relative to clients without those samples, it strictly enlarges the executable experiments on unchanged boundary sorts. The finite-group classification and subgroup-intersection and coupling tests calculate that resource-relative change.

None of these conclusions asserts an old-boundary behavior outside the closure of the entire runtime that performs its acquisition and use, or a general procedure for creating an arbitrary missing physical operation from observations.
\end{theorem}
\begin{proof}
The first claim is Theorems~\ref{m-thm:image} and \ref{m-thm:reuse}. The second is Theorems~\ref{m-thm:sign} and \ref{m-thm:mobius}. The third is Theorem~\ref{thm:frames}, whose inputs do not include a frame. The fourth is Theorem~\ref{thm:frame-reuse} together with the exact identification and later-operation theorems. In the finite implementation, all baseline primitives are equivariant, so the nonconstant relative tests are impossible without the reference. The explicit instructions in Theorem~\ref{g-thm:compiler} realize those tests after acquisition; classification and coupling are Theorems~\ref{g-thm:reference} and \ref{g-thm:coupling}. The final limitation follows from Theorem~\ref{thm:closure}; the geometrical acquisition runtime includes its repetition semantics, and the finite source transfers are expressly accounted for as resources.
\end{proof}

This is not a solution of unrestricted endogenous representation creation. The initial interaction $H$, coefficient model, rank warrant, source samples, and storage are not generated from nothing. In the matrix case the new role is a non-free \emph{type} relative to an old free-role grammar, while its runtime representation and old-boundary computations are conservative. In the finite symmetry case strict new \emph{observed power} comes from additional retained physical samples. The Gaussian bridge constructs rather than merely postulates the frame, but still uses a supplied interaction and a specified sampling runtime. These are distinct gains, now linked by an explicit construction and a shared retained-context theorem rather than conflated.

The stronger problem must specify what it is willing to hold fixed and what can change. Possible nontrivial targets include native fabrication from a specified material and actuator stock, lower bounds on repeated experimental cost, uniform synthesis under unknown-device access, or an evidentially warranted change of primitive operations. This note supplies neither a universal fabrication doctrine nor a completeness theorem across arbitrary native categories. Simply asking for an operation nondefinable in the same closed language that already implements its construction and use cannot supply such a target.

\subsection{Connection to semantic and native doctrines}
The earlier semantic/native distinction is preserved here. A role can be described by its restricted observations without being natively produced. The projector formulas discharge native production in the linear class; the elementary polarity instructions discharge it in a finite class. Conversely, the general $G/H$ theorem is semantic until such an implementation is given. The frame-descent criterion is continuous until its coefficient programs are constructed. These distinctions are explicit hypotheses and conclusions, not fields of a contract left unfilled.

The earlier contextual-reduction and wiring questions also have concrete answers here. An image role is the attainable response object, not automatically a quotient of all previous states. Retaining references changes which tests are possible; the orbit comparison in \eqref{g-eq:coupling} can be an isomorphism, fail injectivity, or fail surjectivity, under the proved finite criteria. Reuse is verified by sandwich, decoder, and intermediate-frame equations rather than by assuming a descent or Beck--Chevalley condition.

The programme's richer contextad setting retains occurrences, endpoint squares, costs, ownership, and proof-relevant native witnesses. Extensional matrix or function equality does not reconstruct those. To lift this construction to that setting one must retain the acquisition and adapter executions and give the required comparison cells and their unit and composition laws. The ordinary coKleisli and slice theorem is a genuine shared categorical statement, but it is only the cartesian retained-resource case; it is not a replacement for a native contextad theorem.

\subsection{Relation to established work}
The categorical and algebraic core is established mathematics. The generated-role category is the Karoubi envelope. Cauchy completion in the ordinary case is completion under retracts of representables \cite[Section 5.8]{Kelly1982}. Finite projectives are summands of finite free modules \cite{Stacks}; with all continuous real coefficients over compact spaces, vector bundles and projective modules are related by Swan's theorem \cite{Swan}. The group classification is the standard orbit form of induction--restriction, and reference stabilizer intersection is elementary group-action theory. The shared slice theorem does not itself constitute a new theory of invention.

Sontag's 1978 split-realization theory is a particularly close comparator \cite{SontagSplit}. It explicitly admits projective, potentially non-free operational state spaces and finite dual-frame reconstruction. For a Noetherian integral domain and a response map of rank $n$ over its fraction field, Theorem 2.1 characterizes split realization by the $n$-minors of the $n$-th block Hankel matrix generating the unit ideal; related reduced-ring statements have their own hypotheses. Our $H$ is a supplied cross-interaction matrix, not automatically a Hankel response operator. We prove no temporal shift invariance, canonical dynamical realization, or minimal-state theorem merely by taking its image. Non-free roles, dual frames, and determinantal realization conditions must therefore not be presented as new ideas here.

Gaussian range acquisition also has direct prior art. Halko, Martinsson, and Tropp describe exact rank-$r$ range recovery from $r$ random probes and the Gaussian range finder $Y=H\Omega$ \cite[Sections 1.3.2 and 4.2, Algorithm 4.1]{HMT2011}. The rank-spanning idea is therefore not new here. The conditioning rejection and elementary probability bound are a transparent refinement within the declared runtime; the role, retained-token, descent, and reuse account is a synthesis whose historical novelty is not established.

For a conservative language extension, Felleisen's Theorem 3.14(i) shows, under his macro-translation conditions, that changing contextual equivalence on old expressions obstructs macro elimination \cite{Felleisen1991}. Our closure and symmetry arguments are direct special comparisons, not replacements for that theorem. They add an explicit resource account and concrete compilers for the classes stated here, not a new general non-eliminability principle.

Konidaris's skill--symbol loop already alternates skill acquisition and representation construction, then acquires skills within the resulting abstract problem \cite{Konidaris2016}. Konidaris, Kaelbling, and Lozano-P\'erez give theorem-backed symbol construction from supplied options, initiation sets, and exact image operators for high-level planning \cite{KKL2018}. Thus the feedback slogan that concepts enable skills and skills enable concepts has direct prior art. Pacheck and collaborators supply learned-representation task encoding and skill repair, while explicitly assuming the needed symbols in their repair setting and leaving physical implementation of suggested skills to the user \cite[Section 6.1 and p.~264]{Pacheck2023}. That is a concrete native-realization gap; it is not closed in general by the narrow frame-acquisition theorem here.

Spivak's polynomial trees already provide supplied interfaces that evolve with interaction, morphisms, tensor structure, and an extension of dynamic organizations \cite{Spivak2026}. The present construction does not establish an equivalence with polynomial trees, infer their objects from arbitrary observations, or supply a native compiler for them. Its contribution to this comparison is the explicit separation of role formation, acquired resources, and implementation warrant. Only its evolving-interface and composition constructions are used for this comparison.

Capucci and Myers give an existing polynomial-monad/contextad transposition theorem \cite[Theorem 4.18]{CM2026}; it must not be renamed as a result of this note. General categorical conceptual blending likewise has existing construction theorems \cite{Blending}. A candidate formed by a descriptive universal property and a resource acquired by a native protocol are different outputs. Whether the precise staged combination proved here is historically new has not been established.

\section{Verification scope}
The theorem and proposition statements establishing the construction include proofs. This source is intentionally delivered uncompiled; no PDF layout verification is claimed.

The separate source \texttt{InventionCore.lean} was checked with Lean 4.35.0 and mathlib revision \texttt{56a005855ea17e0b0958f73ec801ae140b0a1679}. It formalizes the arbitrary-group transitive-token equivalence, its $G/H$ specialization, representative independence, retained composition, the token-elimination criterion, the invariant-source acquisition obstruction, stabilizer intersections, and a strict-subgroup Boolean separator. Its transitive inverse may use classical choice; finite table and polarity implementations above give separate effective realizations. Its reported axioms are only Lean's standard \texttt{propext}, \texttt{Classical.choice}, and \texttt{Quot.sound}, with no unproved placeholders.

That verification does not cover Gaussian acquisition, real arithmetic, topology, finite physical compilers, empirical identification, cost fidelity, contextad witnesses, or historical novelty. The companion finite regression checks cover explicit small action tables and compiler compositions; they are tests, not replacements for the general proofs or physical warrants. The full manuscript has not been formally verified.

\begin{thebibliography}{99}
\bibitem{Felleisen1991}
M. Felleisen, \emph{On the expressive power of programming languages}, Science of Computer Programming 17 (1991), 35--75.
\url{https://www.cs.tufts.edu/comp/150FP/archive/matthias-felleisen/expressive-as-published.pdf}.
\bibitem{Kelly1982}
G. M. Kelly, \emph{Basic Concepts of Enriched Category Theory}, 1982; TAC reprint, Section 5.8.
\url{https://tac.mta.ca/tac/reprints/articles/10/tr10_book_format.pdf}.
\bibitem{Stacks}
The Stacks Project Authors, \emph{Finite projective modules}, Lemma 10.78.2, Tag 00NX.
\url{https://stacks.math.columbia.edu/tag/00NX}.
\bibitem{Swan}
R. G. Swan, \emph{Vector bundles and projective modules}, Transactions of the American Mathematical Society 105 (1962), 264--277.
Bibliographic record: \url{https://www.math.uchicago.edu/~swan/bibliography.html}.
\bibitem{SontagSplit}
E. D. Sontag, \emph{On split realizations of response maps over rings}, Information and Control 37 (1978), 23--33.
\url{https://sontaglab.org/FTPDIR/sontag_minimal_split_realizations_systems_over_rings_information_control1978.pdf}.
\bibitem{Konidaris2016}
G. Konidaris, \emph{Constructing Abstraction Hierarchies Using a Skill-Symbol Loop}, IJCAI 2016.
\url{https://arxiv.org/abs/1509.07582}.
\bibitem{KKL2018}
G. Konidaris, L. P. Kaelbling, and T. Lozano-P\'erez, \emph{From Skills to Symbols: Learning Symbolic Representations for Abstract High-Level Planning}, Journal of Artificial Intelligence Research 61 (2018), 215--289.
\url{https://doi.org/10.1613/jair.5575}.
\bibitem{Pacheck2023}
A. Pacheck, S. James, G. Konidaris, and H. Kress-Gazit, \emph{Automatic encoding and repair of reactive high-level tasks with learned abstract representations}, International Journal of Robotics Research 42(4--5) (2023), 263--288.
\url{https://www.raillab.org/publication/pacheck-2023-automatic/pacheck-2023-automatic.pdf}.
\bibitem{Spivak2026}
D. I. Spivak, \emph{Interactions that reshape the interfaces of the interacting parties}, arXiv:2602.17917v1 (2026).
\url{https://arxiv.org/abs/2602.17917v1}.
\bibitem{CM2026}
M. Capucci and D. J. Myers, \emph{Contextful Arrows via Wreath Products}, arXiv:2410.21889v2, revised 3 November 2026.
\url{https://arxiv.org/abs/2410.21889v2}.
\bibitem{HMT2011}
N. Halko, P.-G. Martinsson, and J. A. Tropp, \emph{Finding Structure with Randomness: Probabilistic Algorithms for Constructing Approximate Matrix Decompositions}, SIAM Review 53(2) (2011), 217--288.
\url{https://amath.colorado.edu/faculty/martinss/2014_CBMS/Biblio/2011_HMT.pdf}.
\bibitem{Blending}
M. Schorlemmer and E. Plaza, \emph{A uniform model of computational conceptual blending}, Cognitive Systems Research 65 (2021), 118--137.
\url{https://www.iiia.csic.es/~enric/papers/UniformModelBlending.pdf}.
\end{thebibliography}
\end{document}
